\chapter{Endoscopy (Upper GI)}

\section*{Overview}
An upper gastrointestinal endoscopy, also called esophagogastroduodenoscopy or EGD, uses a flexible lighted camera to examine the esophagus, stomach, and duodenum. It can obtain biopsies and treat some bleeding or other abnormalities while the clinician is examining the lining.

\section*{Alternative Names}
Esophagogastroduodenoscopy, EGD, Upper endoscopy

\section*{Principle}
A flexible endoscope is passed through the mouth into the esophagus, stomach, and duodenum. Air distends the lumen and the camera gives direct views of the lining, while instruments passed through the scope allow biopsy and treatment.
\section*{History}
Adolf Kussmaul performed the first gastroscopy in 1868 using a rigid tube. Basil Hirschowitz developed the first fully flexible fiberoptic endoscope in 1957.

\section*{Why It Is Done}
Ordered to investigate symptoms of persistent heartburn, nausea, vomiting, difficulty swallowing, abdominal pain, or upper gastrointestinal bleeding.

\section*{Is This a Primary or Secondary Test or Procedure}
Primary diagnostic and therapeutic modality for upper gastrointestinal diseases.

\section*{How It Is Performed}
You lie on your left side. Sedative medication is given IV, and a local anesthetic is sprayed in your throat. The endoscope is gently passed down your esophagus.

\section*{Preparation}
Fast (no food or liquids) for 6 to 8 hours before the procedure. Arrange for a ride home due to sedation.

\section*{Contraindications and Precautions}
Suspected visceral perforation, severe shock, or unstable cervical spine.

\section*{Risks}
Bleeding, infection, or perforation of the GI tract lining (<1 in 1,000 cases). Risks associated with sedation (respiratory depression).

\section*{Aftercare and Recovery}
Rest in recovery until sedation wears off. Do not drive or sign legal documents for 24 hours. Your throat may feel sore for a day.

\section*{Outcome and Follow-up}
The physician inspects the lining for inflammation, ulcers, tumors, or bleeding. Biopsies may be taken for pathological evaluation.

\section*{Expected Results}
Healthy, pink, smooth mucosal lining in the esophagus, stomach, and duodenum; no ulcers, bleeding, or masses.

\section*{Follow-up}
H. pylori testing of biopsy, PPI therapy, or repeat endoscopy.

\section*{When to Seek Medical Care}
Follow the procedure team's specific instructions. Seek urgent medical assessment for severe or worsening pain, uncontrolled bleeding, fever or signs of infection, trouble breathing, chest pain, fainting, new weakness or confusion, inability to urinate, or any rapidly worsening symptom. The appropriate warning signs depend on the procedure and the patient's medical condition.

\section*{References}
\begin{enumerate}
  \item MedlinePlus. Endoscopy (Upper GI). U.S. National Library of Medicine; 2025.
  \item Current Medical Diagnosis and Treatment. McGraw-Hill; 2025.
\end{enumerate}

\clearpage
